\chapter*{源码卷索引}
\phantomsection
\addcontentsline{toc}{chapter}{源码卷索引}

完整源码已经独立成《从 \cpp 到 AI 计算：第三册源码卷》。实践卷不再复制 \filepath{books/cpu-volume-3/labs/linux\_cpu\_inference} 的全部源码，避免一份代码在实践卷和源码卷中重复维护。实践卷负责分阶段安排任务、验收和扩展方向；源码卷负责给出可完整审阅的代码证据。

\begin{sourcereadingbox}
实践卷每一章的代码作业都应回到源码卷对应章节阅读。源码卷导出文件名为 \filepath{从Cpp到AI计算第三册源码卷.pdf} 和 \filepath{从Cpp到AI计算第三册源码卷.epub}，手机目录为 \filepath{/mnt/sdcard/STU/BOOKS/从Cpp到AI计算第三册源码卷}。
\end{sourcereadingbox}

\section*{一、实践卷到源码卷的映射}
\phantomsection
\addcontentsline{toc}{section}{一、实践卷到源码卷的映射}

\begin{description}
  \item[第 1--4 章] 读源码卷第一章和第三章：CMake、README、tiny MLP 模型资产、模型加载、基础推理、int8 kernel 和 CLI。
  \item[第 5--8 章] 读源码卷第二章、第四章和第五章：reference decoder、tokenizer、KV cache trace、sampler、serving probe 和报告字段。
  \item[第 9--12 章] 读源码卷第二章和第四章：safetensors manifest、shape verifier、GPT-2 配置、HuggingFace 权重名映射、memory planner 相关边界。
  \item[第 13--16 章] 读源码卷第三章和第五章：AVX2 路径、benchmark、oneDNN 对标入口、GPT-2 benchmark compare、7B ledger 和 CTest 回归。
\end{description}

\section*{二、最小复现命令}
\phantomsection
\addcontentsline{toc}{section}{二、最小复现命令}

\begin{lstlisting}[language=bash]
cmake -S books/cpu-volume-3 -B books/cpu-volume-3/build/lcqi-debug -DCMAKE_BUILD_TYPE=Debug
cmake --build books/cpu-volume-3/build/lcqi-debug
ctest --test-dir books/cpu-volume-3/build/lcqi-debug --output-on-failure
make cpu3s-pdf
make cpu3s-epub
\end{lstlisting}

完成这组命令后，读者应同时检查实践卷的作业结果和源码卷中的完整实现，确认代码、测试、报告字段和章节讲解是一条证据链。
